\documentclass[a4paper]{report}

% --- FONT & NGÔN NGỮ ---
\usepackage{fontspec}
\setmainfont{Times New Roman}
\usepackage[vietnamese]{babel}
\usepackage{microtype} % Tự động tinh chỉnh khoảng cách để giảm lỗi tràn lề

% --- ÉP FONT 13pt TOÀN CỤC ---
\usepackage{fontsize}
\changefontsize[15.6]{13} % Cỡ chữ 13pt, base leading 15.6pt

% --- CANH LỀ CHUẨN A4 KHÓA LUẬN TỐT NGHIỆP TDTU ---
% Lề trái 3.5cm (đóng gáy sách), lề phải 2.0cm, lề trên 3.0cm, lề dưới 3.0cm
\usepackage[top=3.0cm, bottom=3.0cm, left=3.5cm, right=2.0cm, headheight=16pt, headsep=1.0cm, footskip=1.2cm]{geometry}

% --- DÃN DÒNG & THỤT ĐẦU DÒNG ĐOẠN VĂN ---
\usepackage{setspace}
\onehalfspacing % Dãn dòng 1.5 lines chuẩn
\usepackage{indentfirst}
\setlength{\parindent}{1.27cm} % Thụt đầu dòng 1.27cm (0.5 inch)
\setlength{\parskip}{0pt}      % Không dãn cách đoạn thừa thãi
\emergencystretch 3em          % Chống tuyệt đối lỗi tràn lề phải (overfull hbox)

% --- HEADER & FOOTER: ĐÁNH SỐ TRANG GIỮA TRÊN ---
\usepackage{fancyhdr}
\pagestyle{fancy}
\fancyhf{} 
\fancyhead[C]{\fontsize{10pt}{12pt}\selectfont\thepage} % Số trang căn giữa phía trên
\renewcommand{\headrulewidth}{0pt}
\renewcommand{\footrulewidth}{0pt}

% Đảm bảo trang đầu chương / trang plain vẫn có số trang căn giữa phía trên
\fancypagestyle{plain}{
    \fancyhf{} 
    \fancyhead[C]{\fontsize{10pt}{12pt}\selectfont\thepage} 
    \renewcommand{\headrulewidth}{0pt}
}

% --- TOÁN HỌC & HÌNH ẢNH ---
\usepackage{amsmath, amsfonts, amssymb}
\usepackage{graphicx}
\usepackage{tikz}
\usetikzlibrary{calc, arrows.meta, positioning, fit, backgrounds}
% --- SƠ ĐỒ: bảng màu và biểu tượng dùng chung (tồn kho đa kho) ---
\usetikzlibrary{shapes.geometric, decorations.pathreplacing, shadows}
\definecolor{khoNavy}{HTML}{1F3A5F}
\definecolor{khoCam}{HTML}{E07B24}
\definecolor{khoXanh}{HTML}{2A9D8F}
\definecolor{khoDo}{HTML}{C0392B}
\definecolor{khoNen}{HTML}{F3F5F8}
\tikzset{
    % Cỡ chữ cố định trong sơ đồ, không phụ thuộc cỡ chữ 13pt của thân bài
    soDo/.style={execute at begin picture={%
        \renewcommand\small{\fontsize{10.95}{13.6}\selectfont}%
        \renewcommand\footnotesize{\fontsize{10}{12}\selectfont}%
        \renewcommand\scriptsize{\fontsize{8}{9.5}\selectfont}%
        \renewcommand\tiny{\fontsize{6}{7}\selectfont}}},
    khung/.style={draw=khoNavy!60, fill=khoNen, rounded corners=3pt, line width=0.6pt,
                  align=center, font=\small, inner sep=4pt},
    khungTacTu/.style={khung, draw=khoXanh, fill=khoXanh!10},
    khungDam/.style={khung, draw=khoNavy, fill=khoNavy, text=white},
    khungChiPhi/.style={khung, draw=khoDo!70, fill=khoDo!7},
    nhan/.style={font=\footnotesize, text=khoNavy, align=center},
    muiTen/.style={-{Stealth[length=2.6mm, width=2mm]}, line width=0.9pt, draw=khoNavy},
    muiCam/.style={muiTen, draw=khoCam},
    muiXanh/.style={muiTen, draw=khoXanh},
    muiDo/.style={muiTen, draw=khoDo},
    % Nhà kho: #1 = mức đầy (0..1) của thanh sức chứa
    pics/kho/.style={code={
        \fill[khoNavy!12] (-0.9,0) rectangle (0.9,0.95);
        \draw[khoNavy, line width=0.8pt] (-0.9,0) -- (-0.9,0.95) -- (0,1.45) -- (0.9,0.95) -- (0.9,0) -- cycle;
        \fill[khoNavy!45] (-0.3,0) rectangle (0.3,0.6);
        \foreach \x in {-0.3,-0.1,0.1} \draw[khoNavy!70, line width=0.3pt] (\x+0.1,0) -- (\x+0.1,0.6);
        \foreach \x/\y in {-0.75/0, -0.75/0.22, 0.45/0, 0.45/0.22}
            \filldraw[fill=khoCam!70, draw=khoCam!90!black, line width=0.3pt] (\x,\y) rectangle ++(0.26,0.2);
        \draw[khoNavy, line width=0.5pt] (1.05,0) rectangle (1.25,1.2);
        \fill[khoCam] (1.05,0) rectangle (1.25,{1.2*#1});
    }},
    % Xe tải giao hàng
    pics/xe/.style={code={
        \filldraw[fill=khoCam!80, draw=khoCam!70!black, line width=0.5pt, rounded corners=1pt] (-0.9,0.15) rectangle (0.3,0.85);
        \filldraw[fill=khoNavy!80, draw=khoNavy, line width=0.5pt, rounded corners=1.5pt] (0.32,0.15) -- (0.32,0.65) -- (0.6,0.65) -- (0.85,0.4) -- (0.85,0.15) -- cycle;
        \fill[white!85!khoNavy] (0.4,0.42) rectangle (0.58,0.6);
        \foreach \x in {-0.6,0.55} \filldraw[fill=khoNavy, draw=white, line width=0.5pt] (\x,0.13) circle (0.13);
    }},
    % Khách hàng
    pics/khach/.style={code={
        \foreach \x/\c in {-0.35/khoNavy!55, 0.35/khoNavy!55, 0/khoNavy} {
            \fill[\c] (\x,0.78) circle (0.17);
            \fill[\c] (\x-0.27,0.05) .. controls (\x-0.27,0.55) and (\x+0.27,0.55) .. (\x+0.27,0.05) -- cycle;
        }
    }},
}
\usepackage{float}
\usepackage{placeins}  % \FloatBarrier: giữ bảng, hình của chương trước mục tóm tắt
% Trang chỉ chứa bảng/hình: đặt từ đầu trang thay vì căn giữa theo chiều dọc (tránh khoảng trắng lớn phía trên).
\makeatletter
\setlength{\@fptop}{0pt}
\setlength{\@fpbot}{0pt plus 1fil}
\setlength{\@fpsep}{\floatsep}
\makeatother
% Cho bảng, hình nằm chung trang với chữ nhiều hơn: ít trang chỉ có float, ít khoảng trắng cuối trang.
\renewcommand{\topfraction}{0.9}
\renewcommand{\bottomfraction}{0.8}
\renewcommand{\textfraction}{0.1}
\renewcommand{\floatpagefraction}{0.8}
\usepackage{longtable}
\usepackage{booktabs}
\usepackage{array}
\usepackage{multirow}
\usepackage{tabularx}
\usepackage{xltabular}
\newcolumntype{Y}{>{\centering\arraybackslash}X}
% Bảng kẻ lưới đầy đủ (| và \hline): thêm khoảng đệm trên mỗi hàng để chữ không sát đường kẻ.
\setlength{\extrarowheight}{2pt}
% Ô hai dòng: \tbl, \tbr, \tbc{dòng 1}{dòng 2} căn trái, phải, giữa (tiêu đề cột dài, không phải co bảng).
\newcommand{\tbl}[2]{\begin{tabular}[c]{@{}l@{}}#1\\#2\end{tabular}}
\newcommand{\tbr}[2]{\begin{tabular}[c]{@{}r@{}}#1\\#2\end{tabular}}
\newcommand{\tbc}[2]{\begin{tabular}[c]{@{}c@{}}#1\\#2\end{tabular}}
% \ghichu{...}: ghi chú ngắn dưới bảng (ký hiệu, đơn vị); chú thích bảng chỉ giữ cụm danh từ ngắn.
\newcommand{\ghichu}[1]{\par\smallskip\begin{minipage}{\linewidth}\small\textit{Ghi chú:} #1\end{minipage}}
\usepackage{subcaption}
\graphicspath{{media/}}
\usepackage{caption}
\captionsetup{
    font={small},
    labelfont={bf},
    labelsep=colon,
    skip=6pt
}

\usepackage[hidelinks]{hyperref} % Liên kết mục lục sạch sẽ, không viền đỏ

% --- TÙY CHỈNH MỤC LỤC & DANH MỤC (TOC, LOF, LOT) ---
\usepackage{tocloft} 
\renewcommand{\cftdotsep}{0.5} % Dấu chấm mục lục san sát chuẩn Word

% Thêm chữ "CHƯƠNG " vào Mục lục
\renewcommand{\cftchappresnum}{CHƯƠNG } 
\renewcommand{\cftchapaftersnum}{. }    
\setlength{\cftchapnumwidth}{6.0em}     

% 1. Cấp Chương (Chapter): In đậm toàn hàng
\renewcommand{\cftchapfont}{\bfseries}
\renewcommand{\cftchapleader}{\bfseries\cftdotfill{\cftdotsep}}
\renewcommand{\cftchappagefont}{\bfseries}

\setlength{\cftchapindent}{0pt}

% 2. Cấp Mục (Section - 1.1): Chữ thường, thụt 0.635cm
\renewcommand{\cftsecfont}{\normalfont}
\renewcommand{\cftsecleader}{\normalfont\cftdotfill{\cftdotsep}}
\renewcommand{\cftsecpagefont}{\normalfont}
\setlength{\cftsecindent}{0.635cm}
\setlength{\cftsecnumwidth}{2.2em}

% 3. Cấp Tiểu mục (Subsection - 1.1.1): In nghiêng, thụt 1.27cm
\renewcommand{\cftsubsecfont}{\itshape}
\renewcommand{\cftsubsecleader}{\itshape\cftdotfill{\cftdotsep}}
\renewcommand{\cftsubsecpagefont}{\itshape}
\setlength{\cftsubsecindent}{1.27cm}
\setlength{\cftsubsecnumwidth}{2.8em}

% Dòng đầu của mục, tiểu mục bị mất phần thụt (\parindent của mục lục về 0), nên chèn lại bằng tay.
% Dòng ngắt xuống vẫn thụt theo \leftskip = indent + numwidth nên thẳng hàng với chữ tiêu đề.
\renewcommand{\cftsecfont}{\normalfont\hspace*{\cftsecindent}}
\renewcommand{\cftsubsecfont}{\itshape\hspace*{\cftsubsecindent}}
% Lề phải: số trang rộng 2em, tiêu đề dừng cách lề 3,5em để luôn còn dấu chấm dẫn trước số trang.
\cftsetpnumwidth{2em}
\cftsetrmarg{3.5em}

% Khoảng cách giữa các dòng trong Mục lục
\setlength{\cftbeforechapskip}{8pt} 
\setlength{\cftbeforesecskip}{5pt}
\setlength{\cftbeforesubsecskip}{4pt}

% Danh mục hình vẽ
\renewcommand{\cftfigpresnum}{Hình } 
\renewcommand{\cftfigaftersnum}{: }  
\setlength{\cftfignumwidth}{5.3em}   % đủ cho "Hình 4.10: "

% Danh mục bảng biểu
\renewcommand{\cfttabpresnum}{Bảng }
\renewcommand{\cfttabaftersnum}{: }
\setlength{\cfttabnumwidth}{5.3em}   % đủ cho "Bảng 4.10: "

% Xóa khoảng trắng thừa giữa các chương trong danh mục hình/bảng
\usepackage{etoolbox}
\makeatletter
\AtBeginDocument{
  \patchcmd{\@chapter}{\addtocontents{lof}{\protect\addvspace{10\p@}}}{}{}{}
  \patchcmd{\@chapter}{\addtocontents{lot}{\protect\addvspace{10\p@}}}{}{}{}
}
\makeatother

% --- ĐỊNH DẠNG TIÊU ĐỀ CÁC CẤP (HEADING SPACING CHUẨN ĐẸP) ---
\usepackage{titlesec}
\setcounter{secnumdepth}{3} % Đánh số tới cấp 1.1.1.1

% Heading 1: Chương (\chapter) - 16pt Bold
\titleformat{\chapter}[hang]
  {\fontsize{16pt}{22pt}\selectfont\bfseries} 
  {CHƯƠNG \thechapter.} 
  {0.5cm} 
  {}
\titlespacing*{\chapter}{0pt}{0pt}{16pt} % Cách dưới 16pt

% Heading 2: Mục lớn (\section) - 14pt Bold, sát lề trái 0cm
\titleformat{\section}[hang]
  {\fontsize{14pt}{19pt}\selectfont\bfseries} 
  {\thesection}
  {0.4cm} 
  {}
\titlespacing*{\section}{0pt}{14pt}{6pt} % Sát lề trái

% Heading 3: Tiểu mục (\subsection - vd 3.1.1) - 13pt Bold Italic, thụt vào 0.635cm (nửa tab)
\titleformat{\subsection}[hang]
  {\fontsize{13pt}{18pt}\selectfont\bfseries\itshape} 
  {\thesubsection}
  {0.35cm} 
  {}
\titlespacing*{\subsection}{0.635cm}{10pt}{4pt} % Thụt vào bậc thang 0.635cm

% Heading 4: Tiểu mục con (\subsubsection - vd 3.1.1.1) - 13pt Italic, thụt vào 1.27cm (1 tab)
\titleformat{\subsubsection}[hang]
  {\fontsize{13pt}{18pt}\selectfont\itshape} 
  {\thesubsubsection}
  {0.3cm} 
  {}
\titlespacing*{\subsubsection}{1.27cm}{8pt}{3pt} % Thụt vào bậc thang 1.27cm

% --- KHOẢNG CÁCH CÔNG THỨC TOÁN HỌC ---
\AtBeginDocument{
  \setlength{\abovedisplayskip}{6pt}
  \setlength{\belowdisplayskip}{6pt}
  \setlength{\abovedisplayshortskip}{3pt} 
  \setlength{\belowdisplayshortskip}{3pt}
}

% --- TÀI LIỆU THAM KHẢO ---
\usepackage{csquotes} 
\DeclareQuoteAlias{english}{vietnamese}

\usepackage[backend=biber, style=ieee, language=english, autolang=none, defernumbers=true]{biblatex}
% Tài liệu tham khảo chia theo chủ đề (main.tex in từng nhóm); defernumbers: đánh số liên tục theo thứ tự in.
% Thêm nguồn mới vào references.bib thì phải gán nó vào một nhóm dưới đây, nếu không nguồn đó không được in.
\DeclareBibliographyCategory{tonkho}\addtocategory{tonkho}{harris1913,arrow1951,scarf1960,clark1960,zipkin2000,zipkin2008,silver2016,lee1997}
\DeclareBibliographyCategory{rl}\addtocategory{rl}{bellman1957,puterman1994,sutton2018,schulman2015trpo,schulman2016gae,schulman2017ppo,ng1999,stooke2020,altman1999}
\DeclareBibliographyCategory{marl}\addtocategory{marl}{bernstein2002,oliehoek2016,tan1993,gupta2017,dewitt2020ippo,yu2022mappo,wolpert2001,devlin2011potential,christianos2021seps,zaheer2017deepsets}
\DeclareBibliographyCategory{rltonkho}\addtocategory{rltonkho}{oroojlooyjadid2022,gijsbrechts2022,yang2023mabim,ding2022cdppo,liu2025happo,kotecha2024decentralized,demoor2022reward,zhu2026spatiotemporal,jeong2026alca,boute2022,kaynov2024,vanvuchelen2020,maggiar2025,madeka2022}
\DeclareBibliographyCategory{dubao}\addtocategory{dubao}{ke2017lightgbm,croston1972,syntetos2005accuracy,syntetos2005categorization,kostenko2006,makridakis2022background,makridakis2022accuracy,m5dataset,banrudin2019,elmachtoub2022}
\DeclareBibliographyCategory{danhgia}\addtocategory{danhgia}{brockman2016gym,towers2024gymnasium,hubbs2020orgym,henderson2018,agarwal2021,kunsch1989}
% biblatex khong co vietnamese.lbx -> cac chuoi (vol., no., et al., in Proc.)
% bi in thang ten khoa noi bo (jourvol, andothers). Dung chuoi tieng Anh
% cho danh muc tai lieu, dung quy uoc IEEE.
\DeclareLanguageMapping{vietnamese}{english}
\DefineBibliographyStrings{vietnamese}{%
  url     = {\mkbibbrackets{Online}\adddot\space Available},
  urlseen = {Accessed},
}
\addbibresource{backmatter/references.bib}
\usepackage{xurl}
% Cho phép ngắt DOI/URL sau chữ số và chữ hoa để DOI dài không tràn lề phải.
\setcounter{biburlnumpenalty}{100}
\setcounter{biburlucpenalty}{100}
\let\parencite\cite



\defbibheading{bibcenter}{
    \clearpage
    \phantomsection
    \addcontentsline{toc}{chapter}{TÀI LIỆU THAM KHẢO}
    \begin{center}
        {\fontsize{16pt}{22pt}\selectfont\bfseries TÀI LIỆU THAM KHẢO\par}
    \end{center}
    \vspace{12pt}
}

\defbibheading{subbib}{
    \noindent\textit{#1}\par\vspace{6pt}
}

% Giu moi truong mac dinh cua style IEEE (nhan so [n] treo le trai); chi
% tang khoang cach giua cac muc cho de doc.
\setlength{\bibitemsep}{0.4\baselineskip}



% --- TÙY CHỈNH DANH SÁCH (ITEMIZE / ENUMERATE) ---
\usepackage{enumitem}
\setlist{
    topsep=3pt,
    partopsep=0pt,
    itemsep=2pt,
    parsep=0pt,
    labelindent=1.27cm,
    leftmargin=*,
    align=left,
    labelsep=0.4cm
}
\renewcommand{\labelitemi}{--}
\renewcommand{\labelitemii}{+}
\renewcommand{\labelitemiii}{\textbullet}
\renewcommand{\labelitemiv}{$\cdot$}

\setlist[itemize,1]{
    label=--,
    labelindent=1.27cm,
    leftmargin=*,
    align=left,
    labelsep=0.4cm
}
% Ý con của một gạch đầu dòng dùng dấu "+", nhãn thẳng hàng với chữ của ý cha.
\setlist[itemize,2]{
    label=+,
    labelindent=0pt,
    leftmargin=*,
    align=left,
    labelsep=0.35cm
}
% Ý nhỏ hơn dấu "+" dùng dấu chấm.
\setlist[itemize,3]{
    label=\textbullet,
    labelindent=0pt,
    leftmargin=*,
    align=left,
    labelsep=0.35cm
}
% Gạch đầu dòng nằm trong danh sách đánh số (1), (2)... cũng là ý con nên dùng "+".
\setlist[enumerate]{before={\setlist[itemize,1]{label=+, labelindent=0pt, leftmargin=*, align=left, labelsep=0.35cm}}}